
foudres de la sainte église. Dans cette affaire burlesque,
« le plus âne des trois n'est pas celui qu'on pense »,
aurait dit notre bon La Fontaine. Mais quand on enlève au
peuple le sens du sacré et qu'on pousse la démagogie à le
priver de tout respect, on n'est pas loin d'en arriver à
des actes autrement regrettables. C'est tout le monde qui
en pâtit.

\lettrine{A}{insi}, entre 1555 et 1562, la région de Paulin
est mise en coupe réglée par un certain Cabrol surnommé
« Mandet ». Pour en finir avec ses exactions, on sera obligé
de pendre ce triste personnage. Cabrol mourut sur la potence.

\lettrine{N}{ous} avons vu qu'Etienne de Moret est en 1545
lieutenant du roi pour la baronnie d'Alban ; vers ce même
temps, c'est aussi un de Moret qui est curé de Notre Dame
d'Alban. C'est dire que tant au civil qu'au religieux,
ce sont les « de Moret » qui auront charge de faire face
à l'assaut de la vague protestante qui va déferler dans
notre coin.

\lettrine{E}{st-ce} pour ce motif qu'en août 1565 un messager
venu de Montpellier stationnait à Alban avec son escorte et
se dirigeait ensuite vers Albi pour y rencontrer Dauville,
gouverneur du Languedoc ? ...
